\chapter{Un coucher de soleil}
\label{chap:un-coucher-de-soleil}

Ellundril Chariak était déjà âgée lorsqu'elle rejoignit la caravane.

Très âgée, même.

Pourtant, rien chez elle n'évoquait la fragilité. Elle avait le visage marqué
par les années et de longs cheveux blancs qu'elle portait librement, mais son
corps était resté fin et droit. Ses gestes possédaient une fluidité difficile
à définir. Rien ne semblait jamais forcé chez elle et aucun de ses mouvements
ne paraissait inutile. Lorsqu'elle marchait, contournait un obstacle ou se
baissait pour ramasser quelque chose, chaque geste semblait simplement être
celui qu'il fallait faire.

Plus encore que cette grâce, c'était son harmonie avec ce qui l'entourait qui
frappait Kael. Ellundril ne donnait jamais l'impression de traverser un lieu
qui lui était étranger. Dans une forêt comme au milieu d'un campement, elle
semblait naturellement appartenir à l'endroit où elle se trouvait, adaptant
ses mouvements au terrain plutôt que de chercher à lui imposer les siens.

Kael l'avait remarqué presque immédiatement.

Il avait désormais dix-sept ans et suffisamment appris sur les routes pour
observer la manière dont les gens se déplaçaient. Robin lui avait enseigné à
stabiliser son corps lorsqu'il tirait, Cal à choisir où poser ses pieds et à
utiliser son environnement. Pourtant, Ellundril faisait encore autre chose.

Elle ne semblait jamais lutter contre ce qui l'entourait.

Un terrain irrégulier ne ralentissait pas son pas. Une branche basse ne
l'obligeait pas vraiment à se baisser : son mouvement se modifiait simplement,
avec une souplesse si naturelle qu'il était parfois difficile de déterminer où
finissait sa marche et où commençait l'esquive.

Même Bartiméus semblait avoir accepté sa présence avec une facilité
inhabituelle.

Un soir, peu avant le coucher du soleil, Ellundril vint trouver les deux
garçons alors que la caravane terminait son installation.

— Venez avec moi. J'ai quelque chose à vous montrer.

Kamatsu leva les yeux.

— Où ça ?

Ellundril sourit légèrement.

— Venez.

Kael échangea un regard avec son frère.

Ils la suivirent.

Bartiméus aussi.

\scenebreak

Ellundril les conduisit hors du campement et s'engagea dans la forêt sans
jamais hésiter sur la direction à prendre.

Kael se demanda d'abord si elle connaissait déjà les environs, mais elle
n'avait rien dit qui permette de le savoir. Peut-être était-elle déjà passée
par là. Peut-être avait-elle remarqué quelque chose pendant leur arrivée que
les autres n'avaient pas vu.

Avec elle, il était difficile de savoir.

Le terrain commença bientôt à monter. Ils quittèrent les passages les plus
faciles et progressèrent entre les arbres, les rochers et les racines qui
affleuraient sous la terre. Ellundril avançait devant eux au même rythme que
depuis leur départ.

Kael observa ses pieds.

Elle ne cherchait pas les endroits où marcher. Elle semblait les connaître
avant même d'y poser le pied. Lorsqu'un rocher barrait leur chemin, elle
utilisait son élan pour monter dessus sans ralentir ; lorsqu'une pente devenait
plus raide, elle modifiait légèrement son équilibre et continuait avec la même
aisance. Elle ne combattait jamais le terrain et ne semblait pas davantage
chercher à le dominer. Elle avançait avec lui, utilisant ce qu'il lui offrait
avec une économie de mouvements presque parfaite.

Kael essaya plusieurs fois de reproduire certains de ses gestes.

Il comprit rapidement qu'ils étaient beaucoup moins simples qu'ils en avaient
l'air.

Kamatsu, de son côté, s'intéressait davantage à leur destination.

— On va encore loin ?

— Non.

— Qu'est-ce qu'il y a là-haut ?

Ellundril continua de marcher.

— Tu verras.

Kamatsu regarda Kael.

— Elle répond toujours comme ça ?

— Je crois.

— C'est agaçant.

Ellundril sourit sans se retourner.

Quelques minutes plus tard, les arbres commencèrent à s'espacer.

Puis ils atteignirent le promontoire.

Et Kamatsu cessa de poser des questions.

\scenebreak

La forêt s'étendait sous eux jusqu'aux limites du regard.

Depuis leur hauteur, les arbres n'étaient plus une succession de troncs et de
branches, mais un immense manteau sombre couvrant les vallées et les collines.
Quelques brumes s'attardaient encore dans les creux du relief, dessinant de
longues traînées pâles entre les étendues boisées.

Sur leur droite, la forêt s'ouvrait autour d'un vaste lac. Sa surface presque
immobile s'étirait entre les arbres, découpée par quelques avancées de terre et
de petites îles couvertes de végétation.

Mais ce fut le ciel qui retint leur regard.

À gauche, le soleil touchait l'horizon. La moitié de son disque avait déjà
disparu derrière les terres lointaines et sa lumière embrasait tout ce qui
l'entourait. L'or devenait orange, puis rouge, avant de se mêler progressivement
au rose des nuages.

Les couleurs continuaient à travers toute la voûte céleste.

Le rouge s'adoucissait en rose, le rose se fondait dans le violet et le violet
finissait par céder la place au bleu profond du soir.

Et dans ce bleu, presque à l'opposé du soleil, la pleine lune était déjà levée.

Elle se trouvait plus haut que lui dans le ciel, immense disque clair suspendu
au-dessus de la forêt. Sa lumière demeurait blanche et argentée, mais les
dernières couleurs du crépuscule teintaient encore les nuages qui
l'entouraient. Quelques étoiles commençaient déjà à apparaître dans la partie
la plus sombre du ciel.

Plus bas, le lac capturait une partie de cette lumière. Sa surface reflétait
les violets et les bleus du soir ainsi que la clarté de la lune, tandis que les
forêts qui l'entouraient devenaient progressivement plus sombres.

Kael resta immobile.

Il avait vu des couchers de soleil presque chaque jour depuis qu'il vivait sur
les routes. Il avait vu le ciel devenir rouge au-dessus des plaines, disparaître
derrière des montagnes ou se refléter dans des rivières.

Il n'avait jamais rien vu de semblable.

Ellundril s'était arrêtée entre les deux garçons. Ses longs cheveux blancs
bougeaient doucement dans le vent, mais elle ne disait rien.

Elle n'avait pas besoin de leur expliquer ce qu'ils regardaient.

Kael ne lui demanda pas comment elle avait trouvé cet endroit. Il ne lui
demanda pas davantage comment elle avait su qu'ils devaient y venir
précisément ce soir, à cet instant où le soleil disparaissait tandis que la
pleine lune montait de l'autre côté du ciel.

Peut-être connaissait-elle déjà ce promontoire. Peut-être avait-elle observé la
course du soleil et de la lune pendant les jours précédents. Peut-être avait-elle
simplement remarqué quelque chose que personne d'autre n'avait pris le temps de
voir.

Cela n'avait finalement aucune importance.

Ellundril savait simplement regarder le monde suffisamment bien pour être là
lorsqu'il offrait quelque chose d'exceptionnel.

Kamatsu ne parlait plus non plus.

Pourtant, il aurait probablement pu trouver des dizaines de mots pour décrire
ce qu'il avait devant lui. Depuis Rielle et Thom, il avait appris à observer,
à raconter et à chercher les mots capables de conserver un instant.

Cette fois, aucun ne semblait nécessaire.

Bartiméus s'était installé sur un rocher légèrement à l'écart. Assis bien
droit, la queue ramenée autour des pattes, il contemplait lui aussi le paysage
avec un sérieux inhabituel.

Pendant quelques minutes, personne ne parla.

Le soleil descendit encore.

Les ombres gagnèrent lentement la forêt et les premières étoiles devinrent plus
nombreuses autour de la lune. Sur le lac, les dernières couleurs du jour
s'effacèrent peu à peu au profit du bleu de la nuit.

Le temps semblait s'être arrêté avec eux sur ce promontoire.

Ellundril finit par parler.

— Il y a des choses qu'on comprend mieux lorsqu'on ne cherche pas à les
expliquer.

Ce furent les seuls mots qu'elle prononça.

Kael ne posa aucune question.

\scenebreak

Ils attendirent que le soleil ait entièrement disparu avant de reprendre le
chemin de la caravane.

La forêt était désormais plus sombre, mais Ellundril la traversait avec la même
aisance qu'à l'aller. Kael recommença bientôt à observer sa manière de marcher,
la précision de ses appuis et cette façon qu'elle avait de ne jamais sembler
s'opposer au terrain.

Cette soirée ne fut jamais réellement une leçon.

Ellundril ne leur avait presque rien enseigné et ne leur avait demandé de
retenir quoi que ce soit. Elle les avait simplement conduits jusqu'à un
promontoire au moment où le jour et la nuit semblaient s'y être donné
rendez-vous.

Bien des années plus tard, Kamatsu se souviendrait de son silence et de cette
façon de regarder le monde comme si celui-ci avait toujours quelque chose à
raconter à qui savait lui prêter attention.

Kael, lui, se souviendrait de ses mouvements, de cette fluidité sans effort et
de l'impression qu'Ellundril ne cherchait jamais à vaincre le monde qui
l'entourait, mais avançait avec lui.

Et tous deux se souviendraient du coucher de soleil.